
\section{Installing PostgreSQL}
%% ═════════════════════════════════════════════════════════════════════════════

\begin{frame}[fragile]{\textcolor{brown}{Installing PostgreSQL — Windows}}
  \begin{columns}[T]
    \column{0.55\textwidth}
      \textbf{Installation Steps}
      \begin{enumerate}
        \item Go to \href{https://postgresql.org/download/windows}{\texttt{postgresql.org/download/windows}}
        \item Download the \textbf{Windows x86-64} installer
        \item Run the \texttt{.exe}, grant permissions, click \textit{Next}
        \item Set a password for the default superuser \texttt{postgres}
        \item Keep the default port \textbf{5432} and finish the wizard
      \end{enumerate}

    \column{0.42\textwidth}
      \textbf{Verify \& Connect}
\begin{lstlisting}[style=bash]
psql --version
\end{lstlisting}
\begin{lstlisting}[style=bash]
psql -U postgres
\end{lstlisting}
      \begin{block}{\faKey\; Tip}
        Keep the \texttt{postgres} superuser password safe — you will need it every time you connect.
      \end{block}
  \end{columns}
\end{frame}

% ─────────────────────────────────────────────────────────────────────────────
\begin{frame}[fragile]{\textcolor{brown}{Installing PostgreSQL — Linux}}
  Install PostgreSQL and the contrib extension package:
\begin{lstlisting}[style=bash]
sudo apt update && sudo apt full-upgrade -y
sudo apt install postgresql postgresql-contrib -y
\end{lstlisting}

  Start and enable the service to run on boot:
\begin{lstlisting}[style=bash]
sudo systemctl start postgresql
sudo systemctl enable postgresql
\end{lstlisting}

  Check that the service is running:
\begin{lstlisting}[style=bash]
sudo systemctl status postgresql
\end{lstlisting}

  Switch to the default PostgreSQL user and connect:
\begin{lstlisting}[style=bash]
sudo -i -u postgres
psql
\end{lstlisting}

  You are now inside the \texttt{psql} interactive terminal.
\end{frame}